\scanpage{31}
% L1–4
Supposons donné, pour tte partie finie $A\inserted{\ne\varnothing}$ de $\mathbb N$, un ens $D(A)$ (``drapeau de type $A$'') et pour tte inclusion $A\subset B$, une application ``restriction'' $D(B)\longrightarrow D(A)$, avec \deleted{\unsure{relation|relations}} \inserted{condition} d'identité et de transitivité --
% L5–7
ce qui signifie qu'on a un \deleted{foncteur de la} \inserted{préfaisceau d'ens} sur la cat. \inserted{$C$} associée à l'ens ordonné des parties finies de $\mathbb N$, i.e. $\unsure{\mathbb D|D}\in\operatorname{Ob}\widehat C$.
% L8–9
\deleted{\inserted{\unsure{Posons|Donnons}} Soit d'autre part $\Delta_n=[0,n]\subset\mathbb N$, \inserted{(pour $n\in\mathbb N$)} et soit $\unsure{\alpha|a|\sigma}\colon\Delta_n\hookrightarrow\Delta_m$} Soit
% L10–13
\[
\begin{aligned}
\Theta(\mathbb D)&=\coprod_{\substack{A\subset\mathbb N\\\operatorname{card}A\ \text{fini}}}D\unsure{(|/}A)\\
\Theta_n(\mathbb D)&=\coprod_{\substack{A\unsure{\in|\subset}\mathbb N\\\operatorname{card}A=n+1}}D(A)\qquad n\in\mathbb N
\end{aligned}
\]
% L14
donc
\[
\Theta(D)=\coprod_{n\in\mathbb N}\Theta_n(D)
\]
% L15–17
Dans $\Theta(D)$ on a une relation d'ordre, en posant
\[
\underset{\substack{\rotatebox[origin=c]{90}{$\in$}\\D(A)}}{d}\leq
\underset{\substack{\rotatebox[origin=c]{90}{$\in$}\\D(A')}}{d'}
\overset{\text{déf}}{\Longleftrightarrow}
 A\subset A'\text{ et }d=i_{A,A'}^*(d')
\]
% L18–22
[On peut dire que $D\in\operatorname{Ob}\widehat C$ définit une catégorie fibrée sur $C$ à fibres discrètes, et cette catégorie est associée à un ens ordonné, qui est notre $\Theta(D)$; \unsure{on y a|il y a|on a} une application croissante canonique
\[
\Theta(D)\longrightarrow\unsure{\mathscr P|P|\mathfrak P}_{f}(\mathbb N)
\]
% archive sheet number: 30
